%% kapa: exam=1 round=32 year=2021 period=2 form=A subject=appraisal_law subject_ko=감정평가관계법규 seq=18 number=18
%% source: Q-Net 감정평가사 기출문제 (cst00309) artlSeq=5208397 file=extracted/1008/5208397/A형/2021년도 제32회 감정평가사 1차 2교시 A형.pdf
%% conversion: kordoc 4.24.51 -> tools/gen1.py (auto); review=proofread-sonnet; flags=none
\begin{question}{18}
감정평가 및 감정평가사에 관한 법령상 감정평가사의 권리와 의무에 관한 설명으로 옳지 않은 것은?
\begin{choices}
  \item 등록을 한 감정평가사가 감정평가업을 하려는 경우에는 국토교통부장관에게 감정평가사사무소의 개설신고를 하여야 한다.
  \item 감정평가사는 다른 사람에게 자격증·등록증을 양도·대여하여서는 아니 된다.
  \item 감정평가사는 2명 이상의 감정평가사로 구성된 감정평가사합동사무소를 설치할 수 있다.
  \item 감정평가사는 둘 이상의 감정평가법인 또는 감정평가사사무소에 소속될 수 없다.
  \item 감정평가법인등은 그 직무의 수행을 보조하기 위하여 피성년후견인을 사무직원으로 둘 수 있다.
\end{choices}
\end{question}
\answer{5}
